\chapter{Introdução}
\label{cap:introducao}

O avanço tecnológico, iniciado com as Revoluções Industriais, promoveu a transição da mão de obra manual para a produção mecanizada, resultando em um aumento substancial de produtividade e eficiência. Posteriormente, o advento das \glspl{ict}, como o computador e a Internet, acelerou esse processo ao possibilitar a troca de informação em escala global e viabilizar novos modelos de organização econômica e social. Essas inovações foram determinantes para a intensificação da globalização e da automação, moldando a sociedade contemporânea \cite{Moshawrab2023}. 

Como consequência direta desse desenvolvimento, a ubiquidade de dispositivos computacionais passou a gerar volumes massivos e contínuos de dados. O paradigma tradicional de processamento e armazenamento centralizados torna-se cada vez mais insustentável: além de custos e latência elevados, enfrenta limitações de largura de banda e de escalabilidade, comprometendo a viabilidade de aplicações distribuídas em larga escala, especialmente em \gls{ai} e \gls{ml}. Adicionalmente, arquiteturas centralizadas intensificam preocupações com privacidade e conformidade regulatória, como as impostas pelo \gls{gdpr} \cite{Sikandar2023}.

A \gls{ai} é um campo que estuda algoritmos capazes de executar tarefas tradicionalmente dependentes de inteligência humana. O \gls{ml}, ramo da \gls{ai}, concentra-se em métodos que permitem aos sistemas aprender a partir de exemplos e aprimorar seu desempenho sem programação explícita \cite{Moshawrab2023}.

Nesse contexto, o \gls{fl} surge como uma solução que funciona como um "trabalho em equipe" para modelos de \gls{ai}. Em vez de todos os dados serem enviados para um único lugar, cada participante treina um modelo com seus próprios dados, de forma privada. Depois, eles enviam apenas um resumo do que foi aprendido para um servidor central. A função do servidor é realizar a agregação, que nada mais é do que juntar de forma inteligente o aprendizado de todos para criar um modelo global mais poderoso \cite{Yang2019}. Essa abordagem protege a privacidade, mas cria um desafio: o servidor confia nos resumos que recebe sem nunca ver os dados originais, tornando o sistema vulnerável a ataques durante a fase de agregação \cite{Kairouz2021}.

Essa confiança na etapa de agregação é justamente o ponto fraco do sistema. Um participante mal-intencionado pode enviar um "resumo de aprendizado" falso para contaminar o resultado final, estratégia conhecida como \textit{poisoning attack} (ataque de envenenamento). O ataque pode visar estragar o desempenho geral do modelo (ataque não direcionado) ou criar uma "falha secreta" (\textit{backdoor}), fazendo com que o modelo erre de propósito ao receber um comando específico (ataque direcionado). Usando o método de agregação mais simples, a média, até mesmo um único atacante já consegue causar um dano considerável \cite{Sikandar2023}.

Para mitigar esses ataques, foram propostos agregadores robustos, que atuam como regras aplicadas pelo servidor para combinar as atualizações dos clientes, buscando identificar e reduzir o impacto de contribuições suspeitas. Mecanismos como \textit{Krum}, mediana (\textit{median}) e média aparada (\textit{trimmed mean}) são exemplos de defesas clássicas \cite{Zhang2025}. Entretanto, à medida que novas defesas são desenvolvidas, atacantes também criam estratégias mais sofisticadas para contorná-las, configurando uma dinâmica de coevolução entre mecanismos de defesa e técnicas adversárias \cite{Sikandar2023}.

Uma síntese ampla da área de \gls{fl} discute avanços recentes e reúne problemas e desafios em aberto, reforçando a importância de protocolos de avaliação mais rigorosos e abrangentes ao analisar robustez em cenários adversariais \cite{Kairouz2021}. Em particular, a literatura também discute que escolhas de configuração experimental, como conjunto de dados, distribuição entre clientes e conjunto de ataques avaliados, podem levar a interpretações excessivamente otimistas sobre a robustez de mecanismos de defesa quando o escopo adversarial é limitado \cite{Kairouz2021,Khan2023,Lyu2020}.

Diante desse cenário, a proposta deste trabalho é avaliar a eficácia de diferentes algoritmos de agregação robusta em \gls{fl}. A análise abrange desde abordagens clássicas até propostas recentes baseadas em técnicas de agregação segura e robusta \cite{Zhang2025}, considerando também estratégias que agrupam atualizações por similaridade para lidar com a heterogeneidade de dados \textit{Non-IID} \cite{Zhao2018}. O objetivo é compreender em que condições esses mecanismos são eficazes, suas vulnerabilidades e como cenários experimentais impactam sua avaliação, contribuindo para a consolidação de diretrizes mais consistentes na área de \gls{fl} seguro.
